\subsubsection{Interseccion de segmentos}

\paragraph{Idea intuitiva.}
Dados dos segmentos $[a, b]$ y $[c, d]$, determinar si se intersectan. La idea clasica:
\begin{itemize}\setlength\itemsep{0pt}
	\item Verificar orientaciones: $a, b, c$ y $a, b, d$ deben estar en lados opuestos de la recta $ab$; igual para $cd$ y $ab$.
	\item Caso colineal: si las rectas son paralelas (\texttt{cross} $= 0$) y los segmentos son colineales, verificar que los bounding boxes se intersectan.
\end{itemize}
La demostracion: dos segmentos se intersectan sii cada segmento ``cruza'' la linea que contiene al otro (estricto) o ambos son colineales con solapamiento.

\paragraph{Propiedades clave (invariantes).}
\begin{itemize}\setlength\itemsep{0pt}
	\item Funciona con segmentos en cualquier orientacion.
	\item Caso colineal: requiere check adicional de bounding box.
	\item Para intersectar en un \textbf{punto} especifico: aniadir division con cuidado de precision.
	\item Funciona con segmentos verticales/horizontales sin casos especiales.
\end{itemize}

\paragraph{Codigo clave.}
Test clasico de interseccion:
\begin{verbatim}
int sign(long long x) { return (x > 0) - (x < 0); }
bool intersect(Point a, Point b, Point c, Point d) {
    long long o1 = cross(b - a, c - a);
    long long o2 = cross(b - a, d - a);
    long long o3 = cross(d - c, a - c);
    long long o4 = cross(d - c, b - c);
    if (sign(o1) != sign(o2) && sign(o3) != sign(o4)) return true;
    // casos colineales
    if (o1 == 0 && onSegment(a, b, c)) return true;
    if (o2 == 0 && onSegment(a, b, d)) return true;
    if (o3 == 0 && onSegment(c, d, a)) return true;
    if (o4 == 0 && onSegment(c, d, b)) return true;
    return false;
}
\end{verbatim}

\paragraph{Complejidad.}
\begin{itemize}\setlength\itemsep{0pt}
	\item $O(1)$.
	\item Constante pequena.
\end{itemize}

\paragraph{Donde tocar para modificar.}
\begin{itemize}\setlength\itemsep{0pt}
	\item \textbf{Interseccion con punto}: aniadir la division despues de verificar orientacion.
	\item \textbf{Interseccion rayo-segmento}: ajustar la primera condicion.
	\item \textbf{Poligonos convexos}: usar separacion de ejes (SAT).
	\item \textbf{Multiples queries}: precomputar bounding boxes o usar segment tree.
\end{itemize}

\paragraph{Trampas y errores frecuentes.}
\begin{itemize}\setlength\itemsep{0pt}
	\item Si los segmentos tocan solo en un extremo, el snippet lo considera interseccion (algunos problemas quieren estricto). Ajustar segun el problema.
	\item El caso colineal es comun; olvidar el check adicional de bounding box da falsos negativos.
	\item En problemas con muchos segmentos, considerar sweep line para $O(n \log n)$.
\end{itemize}

\paragraph{Codigo.}
Ver \texttt{cpp.tex}, seccion \textit{Interseccion de segmentos}.

\vspace{0.5em}
